To investigate the effectiveness of an ML model as a first-stage filter for conjunction screening, the primary objective was to find how reliably the model could identify satellite pairs associated with genuine CDM generation using orbital data alone. A particular focus was on the model’s recall rate; a high recall rate means that the model could successfully identify the satellites that should generate a CDM to a high degree of accuracy, therefore minimising the number of potentially dangerous events incorrectly filtered out. 
Historical CDM data was used to provide real-world examples of satellite pairs that had previously flagged a conjunction assessment. Thousands of genuine CDMs were parsed to find pairs of satellites that we know generated a CDM using the current industry method [9]. Historical TLE data from both tracked space objects, which dates back to immediately before the CDM was generated, was then compiled into a positive dataset. In addition to this, a corresponding negative dataset was generated from further TLE data by selecting satellite pairs that satisfied defined non-conjunction criteria and did not produce a CDM under those conditions.
The CDM/no-CDM labels were withheld, leaving only orbital characteristics as the features used by the model in testing. After training the model, a fresh positive and negative testing dataset was then passed through the model. Its predictions were then classified according to whether it was likely to result in CDM generation. These were then compared with the known historical labels, allowing recall and other metrics to be calculated. The model was then evaluated on these metrics to determine its suitability as an initial conjunction screening filter. Outlined in this section are the stages of developing and testing this model.

\subsection{Data Retrieval, Processing, and Labelling}

A large component of the initial programme setup was all about collecting and transforming data. The key data required at this stage were historical CDMs and their time-appropriate TLE data. The necessity for time-appropriate TLE data is crucial. TLE data from after the conjunction assessment may include later observations or manoeuvres that were unavailable at the intended prediction time, potentially producing an unrealistic estimate of the model’s predictive performance.
To retrieve the necessary data (historical CDMs and historical TLEs), an automated session request is made through the Space-Track API [spacetrack website reference] and the specific data is downloaded as a JSON file. This is converted into a CSV file for use with the Pandas data analysis package. Historical data were retrieved in smaller batches to reduce request sizes and avoid timeouts. In addition to this, pauses were added between requests to respect the request-rate limit put in place by Space-Track. Successfully downloaded batches were then saved to a local manifest file, allowing the retrieval to continue after an interruption without repeating completed packets. 
Once all JSON files have been imported and converted appropriately, using the python data analysis library Pandas, the retrieved data was then processed into clean CSV files ready for labelling. In order to produce the positive testing dataset, the retrieved data was converted into one file which contained pairs of satellites with both the CDM data and the first TLE data from immediately before the creation time of the CDM for both satellites. This dataset is all labelled with a ‘1’ indicating that these satellite pairs have resulted in the generation of a CDM. 
Each CDM has a pair key, and this is used to prevent the combination AB being repeated under a new BA check, using up unnecessary computation. These pair keys allow later programmes to directly find key information about specific satellites, including all TLE-present orbital features. A separate programme compiled these features into a single packet associated with each pair, ready to directly pass into an ML model. These packets contained 27 distinct orbital features for every single tracked object. 
Finally, a general TLE pull request was made. A file containing several thousand individual tracked objects and their orbital data was produced. This was processed in a particular way, such that a search of the positive test dataset was made to ensure that the only TLEs left behind were those from tracked objects that have not generated a CDM in the past 3 months. This is the first step to generate our negative test dataset. Each satellite was randomly paired with another from that set. Just because these objects haven’t generated a CDM yet, does not mean that (when paired alongside a second arbitrary object) they wouldn’t produce one in the near future. As a result, these pairs cannot yet be labelled with a ‘0’. This label suggests that no CDM would be generated, which may not be true.

\subsection{Propagation and Strict Negatives}

To avoid this labelling problem for negative pairs, further propagation is required across a 7-day period; this 7-day period is taken directly from industry standards [9].